% Options for packages loaded elsewhere
\PassOptionsToPackage{unicode}{hyperref}
\PassOptionsToPackage{hyphens}{url}
\documentclass[
  brazilian,
]{article}
\usepackage{xcolor}
\usepackage{amsmath,amssymb}
\setcounter{secnumdepth}{-\maxdimen} % remove section numbering
\usepackage{iftex}
\ifPDFTeX
  \usepackage[T1]{fontenc}
  \usepackage[utf8]{inputenc}
  \usepackage{textcomp} % provide euro and other symbols
\else % if luatex or xetex
  \usepackage{unicode-math} % this also loads fontspec
  \defaultfontfeatures{Scale=MatchLowercase}
  \defaultfontfeatures[\rmfamily]{Ligatures=TeX,Scale=1}
\fi
\usepackage{lmodern}
\ifPDFTeX\else
  % xetex/luatex font selection
\fi
% Use upquote if available, for straight quotes in verbatim environments
\IfFileExists{upquote.sty}{\usepackage{upquote}}{}
\IfFileExists{microtype.sty}{% use microtype if available
  \usepackage[]{microtype}
  \UseMicrotypeSet[protrusion]{basicmath} % disable protrusion for tt fonts
}{}
\makeatletter
\@ifundefined{KOMAClassName}{% if non-KOMA class
  \IfFileExists{parskip.sty}{%
    \usepackage{parskip}
  }{% else
    \setlength{\parindent}{0pt}
    \setlength{\parskip}{6pt plus 2pt minus 1pt}}
}{% if KOMA class
  \KOMAoptions{parskip=half}}
\makeatother
\ifLuaTeX
\usepackage[bidi=basic,shorthands=off]{babel}
\else
\usepackage[bidi=default,shorthands=off]{babel}
\fi
\ifLuaTeX
  \usepackage{selnolig} % disable illegal ligatures
\fi
\setlength{\emergencystretch}{3em} % prevent overfull lines
\providecommand{\tightlist}{%
  \setlength{\itemsep}{0pt}\setlength{\parskip}{0pt}}
\usepackage{bookmark}
\IfFileExists{xurl.sty}{\usepackage{xurl}}{} % add URL line breaks if available
\urlstyle{same}
% fallback for those not using the hyperref driver hyperxmp:
\makeatletter
\@ifundefined{xmpquote}{\newcommand{\xmpquote}[1]{#1}}{}
\makeatother
\hypersetup{
  pdftitle={A Ciência do Sagrado e a Estética da Exatidão: Transdisciplinaridade{,} Observação da Natureza e Expressão Artística na Sociedade Maia},
  pdfauthor={DAVID e CARLOS},
  pdflang={pt-BR},
  hidelinks,
  pdfcreator={LaTeX via pandoc}}

\title{A Ciência do Sagrado e a Estética da Exatidão:
Transdisciplinaridade, Observação da Natureza e Expressão Artística na
Sociedade Maia}
\author{DAVID e CARLOS}
\date{Novembro de 2026}

\begin{document}
\maketitle

\section{Introdução}\label{introduuxe7uxe3o}

Em Chichén Itzá, a luz do Sol projeta formas triangulares sobre a
balaustrada norte do Castillo, associadas visualmente ao corpo de uma
serpente. O fenômeno é observável; a intenção de marcar precisamente o
equinócio, porém, não decorre dele de modo automático. Essa diferença
entre o que se registra e o que se infere orienta este ensaio. Em vez de
tomar a arquitetura como prova imediata de uma ciência maia ou como
alegoria de uma integração universal entre arte e conhecimento,
investiga-se o que fontes específicas permitem afirmar sobre a produção
e a comunicação de saberes.

A questão central é como, em contextos maias determinados, observações
do céu, calendários, escrita, pintura e arquitetura organizaram e
comunicaram conhecimentos, e o que esses casos permitem pensar sobre as
divisões modernas entre ciência, arte, religião e política. A hipótese é
que certas práticas e objetos articularam observação empírica,
elaboração calendárica, expressão visual e vida ritual ou política. Essa
articulação é contextual: não pressupõe que os maias nomeassem tais
atividades como disciplinas integradas nos termos atuais, nem que todo
objeto tivesse simultaneamente todas essas funções.

O ensaio compara três conjuntos de evidências: o Castillo de Kukulcán,
em Chichén Itzá, situado por Šprajc e Sánchez Nava no Clássico Terminal
e no Pós-Clássico Inicial; as tabelas astronômicas do Códice de Dresden,
testemunho pós-clássico; e as pinturas de Bonampak, para as quais a
bibliografia disponível é mais limitada. O método distingue dado
material ou documental, interpretação dos pesquisadores, hipótese
debatida e síntese filosófica. A comparação não elimina a desigualdade
do corpus: Chichén Itzá conta com um estudo sistemático de orientações,
enquanto Dresden e Bonampak não dispõem aqui de monografias específicas
suficientes para reconstruir todos os seus procedimentos ou funções.

\section{1. O Marco Teórico: A Fusão Institucional e a Observação da
Natureza}\label{o-marco-teuxf3rico-a-fusuxe3o-institucional-e-a-observauxe7uxe3o-da-natureza}

Johanna Broda propõe estudar os conhecimentos pré-hispânicos em seu
contexto histórico e sociocultural, sem fazer da ciência moderna
ocidental o único parâmetro de racionalidade e sem idealizar realizações
antigas. Isso exige precisão conceitual: ``ciência'', ``natureza'' e
``cosmovisão'' devem ser examinadas em suas manifestações históricas, em
vez de tratadas como categorias abstratas que teriam o mesmo sentido em
todas as sociedades (Broda, 2012, pp.~104--107). A crítica ao
eurocentrismo, portanto, não consiste em chamar qualquer prática antiga
de ciência moderna; consiste em não excluir, por definição, práticas
sistemáticas de observação e previsão porque ocorreram em instituições e
cosmologias diferentes.

Para formular essa abordagem, Broda mobiliza autores de campos
distintos. Ela recorre a Bernal para sublinhar que a ciência tem longa
história e esteve ligada a outras atividades humanas; evoca Einstein
para destacar que os propósitos atribuídos à ciência variam
historicamente; e apresenta uma definição de atividade científica de
Krader, citada por intermédio de Brígida von Mentz, que enfatiza
trabalho empírico, sistemático e cumulativo. Essas referências aparecem
aqui como mediações teóricas apresentadas por Broda, não como consulta
direta às obras de Bernal, Einstein ou Krader. A contribuição central
para este estudo é metodológica: reconhecer práticas empíricas sem
pressupor que tenham constituído uma instituição autônoma equivalente à
ciência contemporânea (Broda, 2012, pp.~106--107).

Broda também retoma a ``polivalência funcional das instituições'',
associada a Terray, e a ``fusão institucional'', associada a Carrasco.
Como exposta por Broda, essa formulação indica que esferas econômicas,
sociais, políticas e ideológicas em sociedades pré-hispânicas não devem
ser isoladas de antemão (Broda, 2012, p.~106). Ela não autoriza concluir
que cada templo fosse ao mesmo tempo observatório, palácio e espaço
ritual, ou que toda imagem tivesse uma função política comprovada.
Trata-se de uma hipótese de trabalho: investigar caso a caso se
evidências diferentes sustentam funções sobrepostas.

Essa cautela se combina com a definição de observação da natureza de
Broda: observação sistemática e repetida dos fenômenos ambientais ao
longo do tempo, capaz de fundamentar previsões e orientar comportamentos
sociais (2012, p.~107). A definição permite reconhecer conhecimento
empírico e cumulativo sem projetar sobre ele uma divisão moderna entre
astronomia, calendário, agricultura e religião. Broda relaciona
observação e cosmovisão, mas não reduz a cosmologia a um disfarce da
política; ambas devem ser entendidas em configurações históricas
concretas (pp.~107--109).

Broda ainda propõe que calendário e observação do horizonte articulavam
tempo e espaço em contextos mesoamericanos (2012, p.~113). Essa
formulação oferece uma pergunta comparativa, não uma conclusão
automaticamente transferível aos maias. Para Chichén Itzá, a ponte é
mais consistente quando apoiada nos dados maias de orientação
arquitetônica analisados por Šprajc e Sánchez Nava. Aveni, em escala
continental, também recomenda considerar conjuntamente orientação,
paisagem e performance ritual (2003, pp.~174--175). Em ambos os casos, a
relação entre astronomia, calendário, prática social e cosmologia deve
ser demonstrada para cada estrutura, e não presumida a partir de um
modelo geral.

\section{2. Estudos de Caso nos Períodos Clássico e
Pós-Clássico}\label{estudos-de-caso-nos-peruxedodos-cluxe1ssico-e-puxf3s-cluxe1ssico}

\subsection{2.1. Templo de Kukulcán em Chichén
Itzá}\label{templo-de-kukulcuxe1n-em-chichuxe9n-itzuxe1}

Šprajc e Sánchez Nava situam Chichén Itzá como centro político
importante no norte das Terras Baixas durante o Clássico Terminal e o
Pós-Clássico Inicial (2013, p.~33). A sequência arquitetônica do sítio é
debatida; Milbrath observa que sua cronologia permanecia controversa
(1988, pp.~57--59). O Castillo deve, por isso, ser tratado dentro de uma
história construtiva complexa, não como produto de um momento único ou
emblema de uma sociedade maia homogênea.

Na escadaria norte, a incidência solar produz triângulos de luz e sombra
sobre a balaustrada. Aveni descreve sete formas luminosas que lembram o
padrão de uma cascavel e observa sua associação visual com as cabeças de
serpente na base da escadaria (2003, p.~163). Essa descrição documenta
um efeito óptico e elementos arquitetônicos que favorecem uma leitura
serpentina. Não demonstra, por si só, que o edifício tenha sido
concebido para assinalar o equinócio ou produzir uma hierofania em uma
data determinada.

A reavaliação de Šprajc e Sánchez Nava situa o Castillo em uma amostra
comparativa de 271 orientações de 87 sítios das Terras Baixas Maias
(2013, p.~35). Os autores observam que o efeito luminoso pode ser
percebido durante vários dias em torno do equinócio e não permite
determinar, apenas pela observação, uma data exata pretendida pelos
construtores (p.~45). Assim, a associação visual com a época equinocial
não basta para provar uma intenção equinocial específica. O resultado
enfraquece uma leitura popular do fenômeno como relógio solar preciso,
mas não elimina a possibilidade de que a arquitetura integrasse
observações sazonais e experiência ritual.

Em escala regional, Šprajc e Sánchez Nava identificam padrões de
orientações solares associados a datas separadas por intervalos
calendáricos e propõem que observações do horizonte pudessem contribuir
para acompanhar o ano sazonal e organizar atividades agrícolas e rituais
(2013, pp.~31--32, 55). Essa interpretação não determina a função
particular do Castillo. Sustenta, com base em recorrências
arquitetônicas, que orientação e calendário podiam relacionar-se a
práticas sociais; a finalidade de cada edifício continua interpretativa.

Milbrath propõe associações entre edifícios e imagens de Chichén Itzá e
fenômenos como o zênite solar, as Plêiades e Vênus (1988, pp.~64--68).
Sua leitura de oito cenas do Templo Superior dos Jaguares como possível
referência ao ciclo de Vênus é uma hipótese iconográfica, não uma
identificação inequívoca (pp.~66--67). A própria autora registra
problemas em interpretações anteriores que combinavam ciclos sem
corresponder a uma situação astronômica real (p.~67). Essas propostas
ampliam as perguntas a investigar, mas não constituem consenso nem
confirmam independentemente a interpretação equinocial do Castillo.

As fontes consideradas não oferecem base suficiente para tratar a
contagem de degraus ou plataformas do Castillo como demonstração de um
calendário de 365 dias. Por isso, a alegação de que a geometria do
monumento codifica um ano solar não é usada como evidência neste ensaio.
A conclusão mais segura é graduada: o efeito de luz e sombra é descrito;
sua integração visual com a escultura serpentina é observável; a relação
com calendário e ritual é plausível no quadro arqueoastronômico
regional; a intenção equinocial exata e uma codificação numérica
específica não estão demonstradas pelo corpus consultado.

\subsection{2.2. Códice de Dresden}\label{cuxf3dice-de-dresden}

O Códice de Dresden é um dos poucos manuscritos maias pré-hispânicos
sobreviventes. Essa preservação limitada torna o códice uma evidência
importante, mas impede tratá-lo como representação completa dos saberes
de todos os povos e períodos maias. Broda destaca o valor dos códices de
Dresden e Madrid para o estudo dos conhecimentos calendáricos e chama-os
de ``verdaderos tratados científicos'' para a pesquisa (2012, p.~111). A
expressão caracteriza sua importância documental; não demonstra que os
produtores organizassem esse saber sob a categoria moderna de ciência.

Milbrath associa ao Dresden o registro do ciclo de Vênus, de
aproximadamente 584 dias, e sua recorrência em um ciclo de oito anos
solares (1988, p.~67). Aveni também menciona a tabela de Vênus do códice
(2003, p.~162), e Broda registra tabelas de lunações e eclipses (2012,
p.~115). Em conjunto, essas referências sustentam que o manuscrito
registra periodicidades celestes em formas tabulares. O corpus
disponível não descreve em detalhe os procedimentos de correção dessas
tabelas nem autoriza afirmar que elas previssem a visibilidade de
eclipses a partir de um lugar determinado.

É importante distinguir esse conteúdo da descrição geral da numeração
maia. Broda discute notação posicional em inscrições e a Conta Longa
(2012, pp.~111, 113); Monteferrante apresenta os sinais de ponto, barra
e concha e explica a estrutura calendárica que inclui o \emph{tun} de
360 dias (2012). Essas informações contextualizam sistemas numéricos
maias, mas não reconstituem os algoritmos das tabelas do Dresden. O
manuscrito não será usado aqui para afirmar uma relação específica entre
o sistema vigesimal, a Conta Longa e cada cálculo astronômico nele
preservado.

O formato do códice permite examinar a disposição de sinais e imagens em
uma superfície organizada para consulta. Houston e Taube argumentam, em
escala mesoamericana, que escrita e dispositivos gráficos podiam evocar
fala e outros sentidos e que a escrita não funcionava apenas como
registro passivo (2000, pp.~261, 265). Aplicada ao Dresden, essa
abordagem orienta uma interpretação sobre a dimensão comunicativa da
página, mas não substitui uma análise codicológica ou iconográfica de
suas seções específicas. A ideia de que sua disposição produz efeito
estético e torna o conhecimento inteligível é uma proposta analítica
deste ensaio, não uma teoria estética nativa comprovada.

O caso demonstra de maneira mais firme a existência de registros
calendáricos e periódicos do que as funções sociais precisas de cada
tabela. A associação entre esses registros, a cosmologia e práticas
rituais é historicamente plausível, mas as fontes aqui reunidas não
permitem reconstruir todos os seus usos. A comparação com a arquitetura
de Chichén Itzá é produtiva porque ambos os casos envolvem regularidades
celestes e formas de registro distintas; não significa que tenham
desempenhado funções equivalentes.

\subsection{2.3. Murais de Bonampak}\label{murais-de-bonampak}

O corpus auditado não contém uma monografia dedicada à cronologia, à
iconografia integral, às técnicas pictóricas ou à história política de
Bonampak. Houston e Taube reproduzem a figura de um cantor e seu glifo
associado aos murais da Sala 1 e discutem a relação entre escrita, fala
e representação (2000, p.~275, fig.~10; pp.~275--276). O artigo situa o
exemplo no contexto da arte maia clássica, mas a referência examinada
não basta para estabelecer a cronologia detalhada do conjunto mural.
Assim, embora o plano de pesquisa proponha situá-lo no Clássico Tardio,
essa datação específica requer uma fonte cronológica própria antes de
ser afirmada com segurança.

O exemplo do cantor sustenta que sinais gráficos podiam participar de
uma representação visual associada à fala. Houston e Taube propõem, de
modo mais amplo, que a percepção e os espaços visuais maias eram
culturalmente organizados, com dimensões hierárquicas e morais (2000,
pp.~261, 287--289). Essa proposta ajuda a formular perguntas sobre a
experiência dos murais, mas não equivale a uma análise de todas as cenas
de Bonampak. Tampouco basta, sozinha, para concluir como cada pintura
foi vista por diferentes públicos ou qual função política específica
desempenhou.

Milbrath relata uma hipótese historiográfica que relaciona imagens de
guerra e simbolismo de Vênus, incluindo paralelos com Bonampak (1988,
pp.~66--67). A autora ressalva que leituras astronômicas não anulam
necessariamente interpretações históricas das cenas (p.~66). Como a
hipótese é secundária e aparece em um estudo centrado na arquitetura e
iconografia astronômica de Chichén Itzá, ela não comprova que o programa
mural de Bonampak seja um calendário, um registro de observações
planetárias ou uma representação de constelações. Não há, nas fontes
auditadas, evidência suficiente para identificar constelações
específicas nas pinturas.

Com a documentação disponível, Bonampak pode ser mobilizado com cautela
para discutir a articulação entre imagem, escrita e representação
narrativa. A relação com memória e legitimação política é uma hipótese
plausível a ser examinada por estudos específicos do sítio, não um
resultado demonstrado pelo exemplo do cantor ou pela referência indireta
a Vênus. O caso complementa o Dresden como outro suporte de comunicação
visual, mas o contraste é assimétrico: para o códice há referências
explícitas a tabelas astronômicas; para os murais, as fontes usadas
sustentam um exemplo gráfico e uma hipótese astral não confirmada.

\section{3. Reflexão Filosófica: A Estética da Exatidão e os Limites do
Conceito Moderno de
Tecnologia}\label{reflexuxe3o-filosuxf3fica-a-estuxe9tica-da-exatiduxe3o-e-os-limites-do-conceito-moderno-de-tecnologia}

``Estética da Exatidão'' é uma categoria criada neste ensaio para
investigar situações em que ordem numérica, repetição calendárica ou
orientação espacial participam de formas de representação e experiência
visual. O conceito não implica que a beleza maia derivasse sempre de
cálculo, nem que os produtores reconhecessem uma unidade entre precisão
e estética nos termos atuais. Sua utilidade depende de distinguir, em
cada caso, medida verificável, regularidade geométrica, organização
calendárica, simbolismo cosmológico, efeito visual e intenção
astronômica --- dimensões que podem coincidir, mas não são sinônimas.

Os três casos oferecem evidências de natureza desigual. No Castillo,
observam-se relações entre forma arquitetônica e luz, enquanto a data e
a intenção do efeito permanecem debatidas. No Dresden, fontes acadêmicas
atestam tabelas de periodicidades celestes, mas o corpus não permite
reconstituir seus procedimentos completos nem todos os usos sociais. Em
Bonampak, as fontes mobilizadas documentam uma figura de cantor e uma
hipótese astral secundária, não um programa astronômico mural. A
comparação sustenta a possibilidade de diferentes suportes participarem
da comunicação do conhecimento, mas não uma síntese segundo a qual todo
edifício, manuscrito e pintura articularia ciência, arte e política da
mesma maneira.

O conceito moderno de tecnologia também requer limites. Se designar
apenas instrumentos voltados a uma finalidade prática, pode ocultar as
dimensões de registro e comunicação presentes em edifícios e
manuscritos; se abranger qualquer objeto simbólico, perde precisão.
Neste ensaio, ``tecnologia'' refere-se, de maneira operacional, a meios
materiais e práticas organizadas de observação, registro e construção,
sem atribuir automaticamente a cada objeto uma função tecnológica ou
científica. Broda e Aveni ajudam a situar conhecimento e prática em
contextos institucionais e rituais; as evidências específicas continuam
a depender de cada caso.

A transdisciplinaridade contemporânea é, portanto, um método de pesquisa
para aproximar arqueologia, astronomia, história, epigrafia e história
da arte. Não é uma continuidade direta de uma disciplina maia nem prova
de que as antigas categorias fossem idênticas às atuais. A contribuição
filosófica dos casos é mais circunscrita: eles mostram que as divisões
modernas são instrumentos analíticos úteis, mas insuficientes quando
naturalizadas como fronteiras universais da experiência histórica.

\section{Conclusão}\label{conclusuxe3o}

A comparação sustenta uma versão delimitada da hipótese da carta de
intenção. Em contextos maias específicos, registros calendáricos,
observações do céu, arquitetura e imagens podiam participar de práticas
de conhecimento e comunicação relacionadas à cosmologia e à vida social.
Broda fornece o eixo teórico para tratar observação empírica e
cosmologia em conjunto, enquanto sua escala mesoamericana exige mediação
por evidências maias específicas. Em Chichén Itzá, o fenômeno luminoso é
observável, mas não fixa a data equinocial nem comprova uma intenção
única. O Dresden documenta tabelas de periodicidades celestes, embora as
fontes reunidas não permitam reconstruir integralmente seus algoritmos
ou usos. Para Bonampak, a documentação aqui disponível é insuficiente
para estabelecer cronologia detalhada, técnica pictórica, função
política específica ou codificação astronômica geral.

Essas diferenças não anulam a pergunta sobre as relações entre arte,
ciência, religião e política; delimitam o que pode ser respondido. Os
casos não demonstram uma fusão universal dessas esferas, mas oferecem
exemplos em que observação, calendário, imagem e espaço construído podem
atravessar fronteiras que a análise moderna tende a separar. A
``Estética da Exatidão'' contribui como lente contemporânea para
examinar tais cruzamentos, desde que não seja projetada como conceito
nativo nem usada para converter qualquer regularidade formal em prova de
intenção científica.

A reflexão transdisciplinar resulta, assim, de um procedimento: partir
dos objetos e manuscritos, identificar o dado disponível, distinguir as
interpretações de seus autores e explicitar o alcance das hipóteses. O
valor contemporâneo do estudo não está em afirmar que os maias
anteciparam nossas disciplinas, mas em mostrar que nossas categorias
precisam ser historicizadas e testadas diante de evidências específicas.

\section{Referências}\label{referuxeancias}

AVENI, Anthony F. Archaeoastronomy in the Ancient Americas.
\emph{Journal of Archaeological Research}, v. 11, n.~2, p.~149--191,
2003.

BRODA, Johanna. Observación de la naturaleza y ``ciencia'' en el México
prehispánico: algunas reflexiones generales y temáticas. In: VON MENTZ,
Brígida (coord.). \emph{La relación hombre-naturaleza: reflexiones desde
distintas perspectivas disciplinarias}. México: CIESAS; Siglo XXI
Editores, 2012. p.~102--135.

HOUSTON, Stephen; TAUBE, Karl. An Archaeology of the Senses: Perception
and Cultural Expression in Ancient Mesoamerica. \emph{Cambridge
Archaeological Journal}, v. 10, n.~2, p.~261--294, 2000. DOI:
10.1017/S095977430000010X.

MILBRATH, Susan. Astronomical Images and Orientations in the
Architecture of Chichen Itza. In: AVENI, Anthony F. (ed.). \emph{New
Directions in American Archaeoastronomy}. Oxford: BAR International
Series 454, 1988. p.~57--82.

MONTEFERRANTE, Sandra. Maya Cycles of Time: Maya Number System.
\emph{Convergence}, Mathematical Association of America, 2012. DOI:
10.4169/loci003886. Disponível em:
\url{https://old.maa.org/node/117895}. Acesso em: 2 nov. 2026.

ŠPRAJC, Ivan; SÁNCHEZ NAVA, Pedro Francisco. Astronomía en la
arquitectura de Chichén Itzá: una reevaluación. \emph{Estudios de
Cultura Maya}, v. 41, p.~31--60, 2013. DOI:
10.19130/iifl.ecm.2013.41.153.

\end{document}
